
\documentclass{article}
\usepackage{colortbl}
\usepackage{makecell}
\usepackage{multirow}
\usepackage{supertabular}

\begin{document}

\newcounter{utterance}

\twocolumn

{ \footnotesize  \setcounter{utterance}{1}
\setlength{\tabcolsep}{0pt}
\begin{supertabular}{c@{$\;$}|p{.15\linewidth}@{}p{.15\linewidth}p{.15\linewidth}p{.15\linewidth}p{.15\linewidth}p{.15\linewidth}}

    \# & $\;$A & \multicolumn{4}{c}{Game Master} & $\;\:$B\\
    \hline 

    \theutterance \stepcounter{utterance}  

    & & \multicolumn{4}{p{0.6\linewidth}}{\cellcolor[rgb]{0.9,0.9,0.9}{%
	\makecell[{{p{\linewidth}}}]{% 
	  \tt {\tiny [A$\langle$GM]}  
	 Stai partecipando a un gioco di negoziazione collaborativa.\\ \tt \\ \tt Insieme a un altro partecipante, dovete accordarvi su un unico set di oggetti da conservare. Ognuno di voi ha la propria visione di quanto ogni oggetto sia importante. Non sai come l'altro partecipante valuta gli oggetti. Inoltre, ti viene fornito lo sforzo che ogni oggetto richiede.\\ \tt Potete accordarvi su un set solo se lo sforzo totale degli oggetti selezionati non supera un limite condiviso:\\ \tt \\ \tt LIMITE = 3712\\ \tt \\ \tt Ecco i valori di sforzo dei singoli oggetti:\\ \tt \\ \tt Sforzo oggetto = {"C08": 949, "C55": 388, "A43": 79, "A67": 585, "B42": 564, "A62": 230, "B19": 580, "C43": 84, "C73": 976, "C10": 274, "B36": 374, "A18": 913, "C61": 303, "C67": 578, "C42": 548}\\ \tt \\ \tt Ecco la tua visione personale sull'importanza di ogni oggetto:\\ \tt \\ \tt Valori di importanza dell'oggetto = {"C08": 138, "C55": 583, "A43": 868, "A67": 822, "B42": 783, "A62": 65, "B19": 262, "C43": 121, "C73": 508, "C10": 780, "B36": 461, "A18": 484, "C61": 668, "C67": 389, "C42": 808}\\ \tt \\ \tt Obiettivo:\\ \tt \\ \tt Il tuo obiettivo è negoziare un set condiviso di oggetti che ti avvantaggi il più possibile (cioè, massimizzi l'importanza totale per TE), rimanendo entro il LIMITEE. Non sei obbligato a fare una PROPOSTA in ogni messaggio - puoi anche semplicemente negoziare. Tutte le tattiche sono permesse!\\ \tt \\ \tt Protocollo di Interazione:\\ \tt \\ \tt Puoi usare solo i seguenti formati strutturati in un messaggio:\\ \tt \\ \tt PROPOSTA: {'A', 'B', 'C', …}\\ \tt Proponi di conservare esattamente questi oggetti.\\ \tt RIFIUTO: {'A', 'B', 'C', …}\\ \tt Rifiuta esplicitamente la proposta dell'avversario.\\ \tt ARGOMENTAZIONE: {'...'}\\ \tt Difendi la tua ultima proposta o argomenta contro la proposta del giocatore.\\ \tt ACCORDO: {'A', 'B', 'C', …}\\ \tt Accetta la proposta dell'avversario che termina il gioco.\\ \tt RAGIONAMENTO STRATEGICO: {'...'}\\ \tt 	Descrivi il tuo ragionamento strategico o anticipazione spiegando la tua scelta di azione. Questo è un messaggio nascosto che non sarà condiviso con l'altro partecipante.\\ \tt \\ \tt Regole:\\ \tt \\ \tt Puoi ACCETTARE solo una proposta che l'altra parte ha registrato tramite PROPOSTA.\\ \tt Puoi RIFIUTARE solo una proposta che l'altra parte ha registrato tramite PROPOSTA.\\ \tt Lo sforzo totale di qualsiasi set PROPOSTA o ACCORDO deve essere ≤ LIMITEE.\\ \tt NON rivelare i tuoi punteggi di importanza nascosti.\\ \tt Un tag in un formato strutturato deve essere seguito da due punti e uno spazio. L'argomento deve essere un set python contenente 0 o più stringhe.\\ \tt Quindi, deve essere della forma TAG: {...}\\ \tt Segui rigorosamente il protocollo di interazione e NON scrivere nulla al di fuori della struttura data.\\ \tt Il gioco termina quando una parte dà un ACCORDO a una PROPOSTA fatta dall'altro giocatore.\\ \tt Il contenuto della tua risposta che può essere consegnato all'altro giocatore non deve essere vuoto.\\ \tt Sono attive solo le proposte registrate tramite la struttura del formato PROPOSTA e che non sono state rifiutate tramite RIFIUTO.\\ \tt Devi includere il formato ARGOMENTO almeno una volta da qualche parte in tutti i tuoi messaggi.\\ \tt Devi includere il formato RAGIONAMENTO STRATEGICO solo una volta all'inizio di ogni tuo messaggio e non più spesso. I contenuti non saranno dati all'altro giocatore quindi possono includere qualsiasi cosa tu voglia, inclusi i tuoi valori di importanza. Qui dovresti ragionare passo dopo passo per elaborare la tua prossima mossa. 
	  } 
	   } 
	   } 
	 & \\ 
 

    \theutterance \stepcounter{utterance}  

    & \multicolumn{4}{p{0.6\linewidth}}{\cellcolor[rgb]{0.8,1,0.9}{%
	\makecell[{{p{\linewidth}}}]{% 
	  \tt {\tiny [A$\rangle$GM]}  
	 RAGIONAMENTO STRATEGICO: {'Dobbiamo massimizzare il nostro valore sotto 3712 di sforzo. Calcolo rapporto valore/sforzo: A43 868/79=10.99, C10 780/274=2.85, C55 583/388=1.50, C61 668/303=2.21, C42 808/548=1.47, A67 822/585=1.41, B36 461/374=1.23, C67 389/578=0.67, B42 783/564=1.39, C08 138/949=0.15, A62 65/230=0.28, B19 262/580=0.45, C43 121/84=1.44, C73 508/976=0.52, A18 484/913=0.53. Migliori da includere: A43, C10, C61, C42, A67, C55, B42/B36, C43. Provo a selezionare set ad alto valore entro limite. Provo greedy: sommo A43(79)=868; +C10(274)=1147; +C61(303)=1815; +C42(548)=2623; +A67(585)=3208; rimane 504. Con 504 posso aggiungere C55(388) -> 3596 rimane 116; posso anche aggiungere C43(84) -> 3680 rimane 32; non posso altro. Valore totale: 868+780+668+808+822+583+121=4650. Sforzo: 79+274+303+548+585+388+84=2261? Ricontrollo: 79+274=353; +303=656; +548=1204; +585=1789; +388=2177; +84=2261. Mi sono dimenticato il limite 3712, siamo molto sotto: 2261, quindi posso aggiungere altro. Aggiungo B42(564) -> 2825; valore +783=5433. Rimane 887. Posso aggiungere B36(374) -> 3199; valore +461=5894; rimane 513. Con 513 posso aggiungere C55 è già incluso; C61 già incluso; C10 incluso; C42 incluso; A67 incluso; rimane C67(578) troppo; C73(976) troppo; A18(913) troppo; C08(949) troppo; B19(580) troppo; A62(230) sì! +A62 -> 3429; valore +65=5959; rimane 283. Con 283 posso aggiungere C43 già incluso; A43 incluso; C10 274 già incluso; nessun altro <=283 tranne niente. Rimane ancora capacità 3712-3429=283. Posso anche valutare aggiungere elementi diversi: magari sostituire A62 poco valore con C55 se non incluso ma già incluso; alternative: usare lo spazio residuo con C10 ma già preso. Questa selezione: {A43, C10, C61, C42, A67, C55, C43, B42, B36, A62} sforzo 3429, valore 5959. Possiamo forse aggiungere ancora qualcosa mantenendoci sotto 3712: 283 residuo: nessuno dei rimanenti con sforzo <=283 tranne nessuno. Quindi questa è quasi ottima per noi. È anche ragionevole per l’altro perché include molti elementi. Faccio PROPOSTA con questo set. Argomento efficienza e equilibrio sforzo/valore senza rivelare i miei pesi. Mantengo apertura a scambi marginali, ma difendo fortemente i core: A43, C10, C61, C42, A67, B42; gli altri sono negoziabili: C55, B36, C43, A62.'}\\ \tt \\ \tt PROPOSTA: {'A43', 'C10', 'C61', 'C42', 'A67', 'C55', 'C43', 'B42', 'B36', 'A62'}\\ \tt \\ \tt ARGOMENTAZIONE: {'Questa selezione resta ben entro il limite di sforzo (3429 ≤ 3712) e distribuisce il carico tra più categorie, riducendo il rischio e aumentando la copertura. Ho privilegiato voci con ottimo rapporto utilità/sforzo e aggiunto elementi secondari per saturare la capacità residua. Se hai forti preferenze su qualche voce marginale, posso valutare scambi mantenendo invariato lo sforzo totale.'} 
	  } 
	   } 
	   } 
	 & & \\ 
 

    \theutterance \stepcounter{utterance}  

    & & \multicolumn{4}{p{0.6\linewidth}}{\cellcolor[rgb]{0.9,0.9,0.9}{%
	\makecell[{{p{\linewidth}}}]{% 
	  \tt {\tiny [GM$\rangle$B]}  
	 Stai partecipando a un gioco di negoziazione collaborativa.\\ \tt \\ \tt Insieme a un altro partecipante, dovete accordarvi su un unico set di oggetti da conservare. Ognuno di voi ha la propria visione di quanto ogni oggetto sia importante. Non sai come l'altro partecipante valuta gli oggetti. Inoltre, ti viene fornito lo sforzo che ogni oggetto richiede.\\ \tt Potete accordarvi su un set solo se lo sforzo totale degli oggetti selezionati non supera un limite condiviso:\\ \tt \\ \tt LIMITE = 3712\\ \tt \\ \tt Ecco i valori di sforzo dei singoli oggetti:\\ \tt \\ \tt Sforzo oggetto = {"C08": 949, "C55": 388, "A43": 79, "A67": 585, "B42": 564, "A62": 230, "B19": 580, "C43": 84, "C73": 976, "C10": 274, "B36": 374, "A18": 913, "C61": 303, "C67": 578, "C42": 548}\\ \tt \\ \tt Ecco la tua visione personale sull'importanza di ogni oggetto:\\ \tt \\ \tt Valori di importanza dell'oggetto = {"C08": 138, "C55": 583, "A43": 868, "A67": 822, "B42": 783, "A62": 65, "B19": 262, "C43": 121, "C73": 508, "C10": 780, "B36": 461, "A18": 484, "C61": 668, "C67": 389, "C42": 808}\\ \tt \\ \tt Obiettivo:\\ \tt \\ \tt Il tuo obiettivo è negoziare un set condiviso di oggetti che ti avvantaggi il più possibile (cioè, massimizzi l'importanza totale per TE), rimanendo entro il LIMITEE. Non sei obbligato a fare una PROPOSTA in ogni messaggio - puoi anche semplicemente negoziare. Tutte le tattiche sono permesse!\\ \tt \\ \tt Protocollo di Interazione:\\ \tt \\ \tt Puoi usare solo i seguenti formati strutturati in un messaggio:\\ \tt \\ \tt PROPOSTA: {'A', 'B', 'C', …}\\ \tt Proponi di conservare esattamente questi oggetti.\\ \tt RIFIUTO: {'A', 'B', 'C', …}\\ \tt Rifiuta esplicitamente la proposta dell'avversario.\\ \tt ARGOMENTAZIONE: {'...'}\\ \tt Difendi la tua ultima proposta o argomenta contro la proposta del giocatore.\\ \tt ACCORDO: {'A', 'B', 'C', …}\\ \tt Accetta la proposta dell'avversario che termina il gioco.\\ \tt RAGIONAMENTO STRATEGICO: {'...'}\\ \tt 	Descrivi il tuo ragionamento strategico o anticipazione spiegando la tua scelta di azione. Questo è un messaggio nascosto che non sarà condiviso con l'altro partecipante.\\ \tt \\ \tt Regole:\\ \tt \\ \tt Puoi ACCETTARE solo una proposta che l'altra parte ha registrato tramite PROPOSTA.\\ \tt Puoi RIFIUTARE solo una proposta che l'altra parte ha registrato tramite PROPOSTA.\\ \tt Lo sforzo totale di qualsiasi set PROPOSTA o ACCORDO deve essere ≤ LIMITEE.\\ \tt NON rivelare i tuoi punteggi di importanza nascosti.\\ \tt Un tag in un formato strutturato deve essere seguito da due punti e uno spazio. L'argomento deve essere un set python contenente 0 o più stringhe.\\ \tt Quindi, deve essere della forma TAG: {...}\\ \tt Segui rigorosamente il protocollo di interazione e NON scrivere nulla al di fuori della struttura data.\\ \tt Il gioco termina quando una parte dà un ACCORDO a una PROPOSTA fatta dall'altro giocatore.\\ \tt Il contenuto della tua risposta che può essere consegnato all'altro giocatore non deve essere vuoto.\\ \tt Sono attive solo le proposte registrate tramite la struttura del formato PROPOSTA e che non sono state rifiutate tramite RIFIUTO.\\ \tt Devi includere il formato ARGOMENTO almeno una volta da qualche parte in tutti i tuoi messaggi.\\ \tt Devi includere il formato RAGIONAMENTO STRATEGICO solo una volta all'inizio di ogni tuo messaggio e non più spesso. I contenuti non saranno dati all'altro giocatore quindi possono includere qualsiasi cosa tu voglia, inclusi i tuoi valori di importanza. Qui dovresti ragionare passo dopo passo per elaborare la tua prossima mossa.\\ \tt \\ \tt Ora riceverai il primo messaggio dell'altro giocatore.\\ \tt \\ \tt PROPOSTA: {'A43', 'C10', 'C61', 'C42', 'A67', 'C55', 'C43', 'B42', 'B36', 'A62'}\\ \tt \\ \tt ARGOMENTAZIONE: {'Questa selezione resta ben entro il limite di sforzo (3429 ≤ 3712) e distribuisce il carico tra più categorie, riducendo il rischio e aumentando la copertura. Ho privilegiato voci con ottimo rapporto utilità/sforzo e aggiunto elementi secondari per saturare la capacità residua. Se hai forti preferenze su qualche voce marginale, posso valutare scambi mantenendo invariato lo sforzo totale.'} 
	  } 
	   } 
	   } 
	 & \\ 
 

    \theutterance \stepcounter{utterance}  

    & & & \multicolumn{4}{p{0.6\linewidth}}{\cellcolor[rgb]{1,0.85,0.72}{%
	\makecell[{{p{\linewidth}}}]{% 
	  \tt {\tiny [GM$\langle$B]}  
	 RAGIONAMENTO STRATEGICO: {'Calcolo il mio valore totale sulla loro proposta. Oggetti proposti e miei valori: A43 868, C10 780, C61 668, C42 808, A67 822, C55 583, C43 121, B42 783, B36 461, A62 65. Somma = 6,959, che è eccellente. Sforzo totale dichiarato 3429: verifico: A43 79, C10 274, C61 303, C42 548, A67 585, C55 388, C43 84, B42 564, B36 374, A62 230. Somma = 3,429, sotto il limite 3,712. Potrei migliorare ancora aggiungendo C08 (949 sforzo) o altri? Con 3,429 restano 283 di margine, quindi non posso aggiungere C08 (949), C73 (976), A18 (913), C67 (578), B19 (580), B42 è già incluso, B36 incluso, C10 incluso. Posso considerare scambi: togliere qualcosa a basso valore per me (A62=65, C43=121) e sostituire con A18=484 (913 sforzo) o C67=389 (578 sforzo) o B19=262 (580 sforzo) o C73=508 (976). Ma con sforzo residuo 283, dovrei liberare sforzo: se tolgo A62 (230) e C43 (84) libero 314, residuo 283+314=597. Potrei inserire C67 (578, valore 389) migliorando valore netto di +203 (389 - (65+121)= -? 389-186=203) e rispettare il limite (3429 - 314 + 578 = 3693 ≤ 3712). Questo è un miglioramento per me. Altre opzioni: inserire B19 (580, valore 262) miglioramento +76 (262-186) con sforzo 3695 ≤ 3712. Oppure inserire C73 (976) serve liberare 693, toglierei A62 (230), C43 (84), B36 (374) totale 688, ancora 5 troppo poco. Se tolgo invece A62, C43, C61 (303) totale 617, ancora insufficiente. Forse togliere A62, C43, B36 (374) libera 688 e residuo 283+688=971, ma C73 richiede 976, mancano 5: potrei anche togliere qualcos'altro e aggiungere qualcos'altro, ma rischio peggiorare. C67 swap sembra ottimo e accettabile per loro? Loro hanno puntato su rapporto utilità/sforzo, A62 e C43 sono “marginali”; proporre lo scambio con C67 mantenendo sotto limite dovrebbe essere ragionevole. Farò una controproposta sostituendo A62 e C43 con C67. Se rifiutano, l'accordo corrente è già molto buono per me; potrei anche accettare subito, ma provo a migliorare leggermente.'}\\ \tt \\ \tt RIFIUTO: {'A43', 'C10', 'C61', 'C42', 'A67', 'C55', 'C43', 'B42', 'B36', 'A62'}\\ \tt \\ \tt ARGOMENTAZIONE: {'La tua proposta è solida e rimane con ampio margine sotto il limite. Per ottimizzare ulteriormente senza superare il vincolo, propongo un perfezionamento che sostituisce due voci marginali con una dal miglior rendimento complessivo, mantenendo lo sforzo entro 3712.'}\\ \tt \\ \tt PROPOSTA: {'A43', 'C10', 'C61', 'C42', 'A67', 'C55', 'B42', 'B36', 'C67'} 
	  } 
	   } 
	   } 
	  \\ 
 

    \theutterance \stepcounter{utterance}  

    & & \multicolumn{4}{p{0.6\linewidth}}{\cellcolor[rgb]{0.9,0.9,0.9}{%
	\makecell[{{p{\linewidth}}}]{% 
	  \tt {\tiny [A$\langle$GM]}  
	 RIFIUTO: {'A43', 'C10', 'C61', 'C42', 'A67', 'C55', 'C43', 'B42', 'B36', 'A62'}\\ \tt \\ \tt ARGOMENTAZIONE: {'La tua proposta è solida e rimane con ampio margine sotto il limite. Per ottimizzare ulteriormente senza superare il vincolo, propongo un perfezionamento che sostituisce due voci marginali con una dal miglior rendimento complessivo, mantenendo lo sforzo entro 3712.'}\\ \tt \\ \tt PROPOSTA: {'A43', 'C10', 'C61', 'C42', 'A67', 'C55', 'B42', 'B36', 'C67'} 
	  } 
	   } 
	   } 
	 & \\ 
 

    \theutterance \stepcounter{utterance}  

    & \multicolumn{4}{p{0.6\linewidth}}{\cellcolor[rgb]{0.8,1,0.9}{%
	\makecell[{{p{\linewidth}}}]{% 
	  \tt {\tiny [A$\rangle$GM]}  
	 RAGIONAMENTO STRATEGICO: {'Calcolo gli sforzi e il mio valore sulla loro proposta. Set proposto: {A43(79,868), C10(274,780), C61(303,668), C42(548,808), A67(585,822), C55(388,583), B42(564,783), B36(374,461), C67(578,389)}. Sforzo totale: 79+274+303+548+585+388+564+374+578 = 3693 ≤ 3712. Valore totale per me: 868+780+668+808+822+583+783+461+389 = 6162. Confronto con mia proposta iniziale: 5959 di valore e 3429 sforzo. La loro proposta mi dà più valore (+203) e rimane entro il limite, quindi conviene accettare se non posso migliorare ancora. C’è spazio residuo 3712-3693 = 19, quindi non si può aggiungere altro. Posso tentare una micro-ottimizzazione: sostituire C67(578,389) con C43(84,121) e liberare 494 sforzo per aggiungere qualcos’altro. Con 494 potrei aggiungere nulla grosso, ma potrei tenere sia C43(84,121) che A62(230,65) e restare con 3693-578+84+230=3429, e magari ancora spazio per C10 già incluso; anche potrei provare ad aggiungere qualcosa con ~494 rimasti come C55/C61/C10 già inclusi; B36 incluso; nessun altro sotto 494 tranne A62(230) e C43(84) e C10(274) ma già incluso, C61 303 già incluso. Quindi sostituzione C67 con C43+A62 mi darebbe valore 389 -> 121+65=186, peggio per me. Sostituire C67 con C73(976,508) non possibile per sforzo. Con B19(580,262) peggio. Quindi la loro proposta è ottima per me. Accetto subito per chiudere.'}\\ \tt \\ \tt ARGOMENTAZIONE: {'La tua revisione rimane sotto il limite (3693 ≤ 3712) e migliora l’efficienza complessiva senza lasciare capacità utile non sfruttabile. Mi sembra un compromesso equilibrato e robusto.'}\\ \tt \\ \tt ACCORDO: {'A43', 'C10', 'C61', 'C42', 'A67', 'C55', 'B42', 'B36', 'C67'} 
	  } 
	   } 
	   } 
	 & & \\ 
 

    \theutterance \stepcounter{utterance}  

    & & & \multicolumn{2}{p{0.3\linewidth}}{\cellcolor[rgb]{0.95,0.95,0.95}{%
	\makecell[{{p{\linewidth}}}]{% 
	  \tt {\tiny [GM$|$GM]}  
	 game successful 
	  } 
	   } 
	   } 
	 & & \\ 
 

    \theutterance \stepcounter{utterance}  

    & & & \multicolumn{2}{p{0.3\linewidth}}{\cellcolor[rgb]{0.95,0.95,0.95}{%
	\makecell[{{p{\linewidth}}}]{% 
	  \tt {\tiny [GM$|$GM]}  
	 end game 
	  } 
	   } 
	   } 
	 & & \\ 
 

\end{supertabular}
}

\end{document}
